\section{Naturaleza de los datos}

Hay dos grupos de datos para el proyecto, uno perteneciente al “Problema de colas” y otro para la obtención del BCS por medio de imágenes. Los datos del problema de colas los generan los robots de ordeño DeLaval VMS del establo en CAETEC. Estos datos fueron compartidos por el profesor Benjamín Valdéz Aguirre y el profesor Eduardo Daniel Juárez Pineda por medio de una carpeta compartida en Google Drive. Para el dataset de fotos, las fotos fueron tomadas por estudiantes de los equipos del reto el 29 de septiembre de 2026, con dispositivos iPhone.

\begin{table}[htbp]
  \centering
  \small
  \begin{tabularx}{\textwidth}{>{\raggedright\arraybackslash}p{0.18\textwidth}X>{\raggedright\arraybackslash}p{0.16\textwidth}>{\raggedright\arraybackslash}p{0.13\textwidth}}
    \toprule
    Conjunto & Contenido & Formato & Periodo \\
    \midrule
    Ordeños por vaca & Hora, duración, producción, flujos por cuarto de ubre, conductividad, sangre, RCS, patadas, incompletos & 36 CSV, uno por vaca & mar–jun 2025 \\
    Reporte reproductivo & Estado reproductivo y días desde parto, celo e inseminación & CSV & jul 2025 \\
    Inventario & Grupo, edad, concentrado consumido, fecha esperada de parto & CSV & jul 2025 \\
    Patadas & Registro de patadas por vaca & CSV & jul 2025 \\
    Histórico & Producción por sesión (vaca 6229) y paquete “Datos 2026 TEC-UAQ-UNAM” & XLSX, ZIP & hasta ene 2026 \\
    Condición corporal & BCS asignado por expertos & Hoja de Google & N/A(?) \\
    Fotos & Imágenes de vacas, una carpeta por ID & JPG & sep 2026 \\
    \bottomrule
  \end{tabularx}
\end{table}
